\documentclass[11pt]{article}
\usepackage[margin=2.1cm]{geometry}
\usepackage{booktabs}
\usepackage{amsmath}
\usepackage{enumitem}
\usepackage[T1]{fontenc}
\usepackage{lmodern}

\setlength{\parindent}{0pt}
\setlength{\parskip}{5pt}
\newcommand{\code}[1]{\texttt{#1}}
\newcommand{\answerbox}[1]{\vspace{2pt}\fbox{\parbox[c][#1][c]{\dimexpr\linewidth-2\fboxsep-2\fboxrule}{\ }}\vspace{4pt}}

\begin{document}

\begin{center}
{\Large\bfseries EE\,362 --- Pre-Lab 5}\\[2pt]
{\large Arithmetic Logic Unit}\\[6pt]
\begin{tabular}{@{}p{6.5cm}p{4cm}p{4cm}@{}}
Name: \hrulefill & Section: Mon / Tue & Date: \hrulefill
\end{tabular}
\end{center}

\vspace{2pt}\hrule\vspace{4pt}

{\small\textbf{Due at the start of the lab session.} Five points, about twenty
minutes, on paper. This one is worth doing properly --- the ALU is the biggest
component so far and the students who arrive with these answers finish.}

\vspace{4pt}\hrule

\section*{Question 1 --- Paper first}

\textbf{(a)} Fill in the 32-bit result, in hex, for each operation. Assume
\code{Shft = 1} where a shift amount is needed.

\begin{center}
\begin{tabular}{llll@{\hskip 8pt}||@{\hskip 8pt}l}
\toprule
\# & operation & \code{Ain} & \code{Bin} & \code{Rout} \\
\midrule
1 & \code{SUB}  & \code{0x00000004} & \code{0x00000007} & \rule{3.6cm}{0.4pt} \\[5pt]
2 & \code{NOT}  & \code{0x00000000} & \code{0x0F0F0F0F} & \rule{3.6cm}{0.4pt} \\[5pt]
3 & \code{SRL}  & \code{0xFFFFFFF8} & --- & \rule{3.6cm}{0.4pt} \\[5pt]
4 & \code{SRA}  & \code{0xFFFFFFF8} & --- & \rule{3.6cm}{0.4pt} \\[5pt]
5 & \code{SOLS} & \code{0xFFFFFFFF} & \code{0x00000001} & \rule{3.6cm}{0.4pt} \\[5pt]
6 & \code{SOLU} & \code{0xFFFFFFFF} & \code{0x00000001} & \rule{3.6cm}{0.4pt} \\
\bottomrule
\end{tabular}
\end{center}

\textbf{(b)} Rows 3 and 4 have the same input and different answers. Rows 5 and 6
also. In one sentence each, say what information decides the correct answer.

\answerbox{2.0cm}

\textbf{(c)} Write the two flag equations for a 32-bit result $R$:

\vspace{4pt}
$\text{Zero} = $ \rule{9cm}{0.4pt}
\vspace{8pt}

$\text{Neg} = $ \rule{9cm}{0.4pt}
\vspace{6pt}

One sentence: why does \code{Neg} require no computation at all?

\answerbox{1.2cm}

\section*{Question 2 --- Interface and structure}

\textbf{(a)} Write the VHDL \code{entity} for \code{ALU}. Interface only.

\answerbox{4.2cm}

\textbf{(b)} An ALU computes every operation in parallel and multiplexes the
answers. Suppose you took the \code{Zero} flag straight from the adder instead
of from the selected result. Give one specific operation and operand pair that
would then report the wrong flag, and say what the flag would be versus what it
should be.

\answerbox{2.2cm}

\section*{Question 3 --- Test plan}

You cannot test $2^{64}$ operand pairs. State your coverage argument, then give
the vectors.

\begin{center}
\begin{tabular}{@{}c@{\hskip 8pt}p{4.6cm}@{\hskip 8pt}p{8.0cm}@{}}
\toprule
\# & operation and operands & which wrong implementation this catches \\
\midrule
1 & \rule{4.4cm}{0.4pt} & \rule{7.8cm}{0.4pt} \\[8pt]
2 & \rule{4.4cm}{0.4pt} & \rule{7.8cm}{0.4pt} \\[8pt]
3 & \rule{4.4cm}{0.4pt} & \rule{7.8cm}{0.4pt} \\[8pt]
4 & \rule{4.4cm}{0.4pt} & \rule{7.8cm}{0.4pt} \\[8pt]
5 & \rule{4.4cm}{0.4pt} & \rule{7.8cm}{0.4pt} \\[8pt]
6 & \rule{4.4cm}{0.4pt} & \rule{7.8cm}{0.4pt} \\
\bottomrule
\end{tabular}
\end{center}

\textbf{Then answer both, one or two sentences each:}

\textbf{(a)} Which single vector in your table would catch an \code{SRA} that was
accidentally written as an \code{SRL}? If you do not have one, add it now.

\answerbox{1.6cm}

\textbf{(b)} For arithmetic, ``test the boundaries'' replaces ``test
exhaustively''. Name the five boundary values worth testing on a 32-bit signed
operand, and say why the middle of the range is uninteresting.

\answerbox{2.2cm}

\vfill
{\small\emph{Question 3 is the one that matters.} It is the same question as
Section 5 of your lab report.}

\end{document}
